\par\Needspace{9\baselineskip}\bigskip
\pdfbookmark[3]{Exposition, 19}{pallas-II-article19}
\begin{center}19.\end{center}
\begingroup
\setlength{\lineskip}{4pt}\setlength{\lineskiplimit}{4pt}
\noindent\hspace*{\parindent}Applying these methods of integration to our variations in the elements gives, apart from the constants introduced by integration, the values of \(i,\ \AscendingNode,\ \varphi,\ n,\ \tilde\omega,\ \epsilon\) at epochs intermediate to those for which the instantaneous variations were calculated, and the value of \(\int n\mathrm{d}t\) at the original epochs, if formulas I, II are used, or vice versa if formulas III, IV are preferred. One can also use formula III for single integration and formula II for double integration. This gives all the elements at the same epochs. The choice is of little consequence, since these values of the elements need serve only to derive, by the known methods of interpolation, the values corresponding to the epochs of the different observations. The calculations from these observations will yield, by known methods, the values of the constants. There are seven of them, but they reduce to six because the constant \(\alpha\) resulting from \(\int\dfrac{\mathrm{d}\epsilon}{\mathrm{d}t}\mathrm{d}t\) and the constant \(\beta\) arising from \(\int n\mathrm{d}t\) or \(\int\!\int\dfrac{\mathrm{d}n}{\mathrm{d}t}\mathrm{d}t^2\) combine into a single constant.
\par\endgroup
